% ═══════════════════════════════════════════════════════════════════════════════
%  CAPÍTULO — Abstract Factory
% ═══════════════════════════════════════════════════════════════════════════════

\chapter{Abstract Factory}

\patroninfo{Creacional}{87}

% ─── Situación Problema ──────────────────────────────────────────────────────
\section{Situación Problema}

Airbnb opera en tres superficies distintas: aplicación web, aplicación
iOS y aplicación Android. Cada superficie tiene sus propios estándares
visuales: los botones, tarjetas de propiedad y formularios de reserva
difieren en apariencia y comportamiento entre plataformas, pero siempre
deben mantenerse visualmente coherentes \emph{dentro} de la misma
plataforma.

El equipo de frontend se encontró con que el código de composición de
pantallas mezclaba condicionales de plataforma con la lógica de negocio.
Crear un componente de tarjeta de propiedad para web y luego reutilizar
ese mismo constructor en la app de iOS producía inconsistencias visuales.
El problema se agravaba cada vez que se quería incorporar una nueva
plataforma: había que rastrear todos los puntos del código donde se
instanciaban componentes y añadir una nueva rama condicional.


% ─── Solución ────────────────────────────────────────────────────────────────
\section{Solución}

El patrón Abstract Factory define una interfaz \texttt{FabricaUI} con
métodos para crear cada familia de componentes (\texttt{crearBoton()},
\texttt{crearTarjetaPropiedad()}, \texttt{crearFormularioReserva()}).
Las fábricas concretas \texttt{FabricaWeb}, \texttt{FabricaIOS} y
\texttt{FabricaAndroid} implementan esa interfaz produciendo componentes
coherentes con su plataforma.

El código de composición de pantallas trabaja únicamente contra
\texttt{FabricaUI}; en tiempo de ejecución se inyecta la fábrica
correspondiente a la plataforma activa. Esto garantiza que los componentes
de una misma familia siempre sean compatibles entre sí, elimina los
condicionales de plataforma dispersos por el código y permite incorporar
nuevas superficies simplemente añadiendo una nueva fábrica concreta.

% ─── Diagrama de clases ─────────────────────────────────────────────────────
\begin{figure}[H]
    \centering
    % \includegraphics[width=\textwidth]{imagenes/abstract_factory_diagrama.png}
    \caption{Diagrama de clases del patrón Abstract Factory}
\end{figure}


% ─── Roles de las Clases ────────────────────────────────────────────────────
\section{Roles de las Clases}

% Explicar la responsabilidad de cada clase/interfaz

\begin{itemize}
    \item \textbf{AbstractFactory:}
    \item \textbf{ConcreteFactory:}
    \item \textbf{AbstractProduct:}
    \item \textbf{ConcreteProduct:}
    \item \textbf{Client:}
\end{itemize}


% ─── Main ───────────────────────────────────────────────────────────────────
\section{Main}

% Explicar qué realiza el programa principal

\begin{lstlisting}[language=Java, caption=Ejemplo de clase Java]
public class Usuario {

    private String nombre;
    private int edad;

    public Usuario(String nombre, int edad) {
        this.nombre = nombre;
        this.edad = edad;
    }

    public void mostrarInformacion() {
        System.out.println("Nombre: " + nombre);
        System.out.println("Edad: " + edad);
    }

    public static void main(String[] args) {

        Usuario usuario = new Usuario("Carlos", 25);

        usuario.mostrarInformacion();
    }
}
\end{lstlisting}


% ─── Salida del Main ────────────────────────────────────────────────────────
\section{Salida del Main}

\begin{lstlisting}[language=Java, caption=NombreClase]
 
\end{lstlisting}


% ─── Clases ─────────────────────────────────────────────────────────────────
\section{Clases}

% ─── Clase 1 ────────────────────────────────────────────────────────────────
\subsection{NombreClase}

\begin{lstlisting}[language=Java, caption=NombreClase]
 
\end{lstlisting}


% ─── Clase 2 ────────────────────────────────────────────────────────────────
\subsection{NombreClase}

\begin{lstlisting}[language=Java, caption=NombreClase]
 
\end{lstlisting}


% ─── Clase 3 ────────────────────────────────────────────────────────────────
\subsection{NombreClase}

\begin{lstlisting}[language=Java, caption=NombreClase]
 
\end{lstlisting}


% ─── Conclusiones ───────────────────────────────────────────────────────────
\section{Conclusiones}

% Explicar:
% - Ventajas del patrón
% - Cuándo usarlo
% - Beneficios obtenidos
% - Posibles desventajas
% - Relación con principios SOLID